% a06_BALtorus_d1: part of Step 3 of the proof of thm:a06:T1 (input from sec/a06_BALtorus_c.tex).
% Group (B): the perturbative form of the effective action, Ba1 (3.34)-(3.36).

\subsubsection*{Notation for Steps 3(B)--3(D)}
Throughout the rest of Step 3, $d=3$, $L=5$, $k$ is the step, $t_k\defeq L^k\eps\le\eps_0\le1$ is Ba{\l}aban's scale parameter (written $t$ when $k$ is fixed), $\eta\defeq L^{-k}$, and all scale-$k$ operators act on the unit lattice $T^{(k)}_1$ of scale $k$ with the fine lattice $\eta\Z^3$ inside each unit cell. The rescaled couplings are $e_k\defeq e\,t_k^{1/2}$, $\lambda_k\defeq\lambda t_k$; $p(t)\defeq b_0(1+\ln t^{-1})^{\bar p}$ and $r(t)\defeq R(1+\ln t^{-1})^{\varrho}$, $\varrho>1$, are the field bound and the layer width of \cite[(2.7), (2.16)--(2.17)]{Ba2}; $\Lambda^{(k)}_i$, $i=-1,\dots,8$, are the regions of \cite[(2.8)]{Ba2} at scale $k$ (unions of large blocks, $\Lambda^{(k)c}_{i+1}$ the union of the large blocks within distance $r(t_k)$ of $\Lambda^{(k)c}_i$), and a primed region $\Lambda^{(k-1)'}_i$ is the scale-$(k-1)$ region read on $T^{(k)}_1$. $M$ is the side of Ba{\l}aban's large blocks (\cite[p.~607]{Ba1}: the block partition with $L$ replaced by $ML$; written $M_b$ in \cref{thm:a06:BALA}), not the scalar mass of \cref{thm:a06:T1}, which is $m$ or $M_1$ here. The expansion order is $\bar n$; $a_k$ is the block-spin coefficient of \cite[(2.13)]{Ba1}, with $a_{k+1}=aa_k/(aL^{-2}+a_k)$, $a_1=a$, so that $a_*\defeq a(1-L^{-2})\le a_k\le a$ for all $k$. $Q_k(A)$ is the covariant $k$-step block average, $P_k=Q_k^*Q_k$, $G_k(\Omega,A)=(-\Delta^{\eta,N}_{A,\Omega}+m_k^2+a_kP_k(A))^{-1}$ the scalar propagator on a union $\Omega$ of blocks with Neumann conditions, $\Delta^{(k)}(\Omega,A)=a_k-a_k^2Q_kG_k(\Omega,A)Q_k^*$ the block-spin operator, $C^{(k)}_\Lambda=((\Delta^{(k)}+aL^{-2}P)|_\Lambda)^{-1}$ the fluctuation covariance, as in \cite[Sect.~2]{Ba1} and \cite[Sect.~1]{Ba4}. For a kernel or external field, $\norm{\cdot}_{1,\alpha_0}$ is the norm \cite[(1.32)]{Ba3} (sup, one covariant derivative, its H\"older quotient of exponent $\alpha_0$), in each variable for kernels of two variables. ``Proposition~1'' is \cite[Prop.~1, (1.33)]{Ba3}: for a localised renormalised class $G_{\mathrm{ren}}$ with localisation $\Lambda=\{\square(v)\}$, $|E(G_{\mathrm{ren}},\Lambda,\Phi)|\le O(1)e_k^{d_v}\lambda_k^{d_s}e^{-\delta_0d(\Lambda)}\prod\norm{h\Phi}_{1,\alpha_0}$, $d(\Lambda)$ the length of a shortest tree through the cubes of $\Lambda$; its proof at $L=5$ is Step~3(D) below.

\subsubsection*{Step 3(B): the perturbative form of the effective action, \cite[(3.34)--(3.36)]{Ba1}}

\cite[(3.36)]{Ba1} asserts, ``by close inspection'', that after $k$ steps the polynomial part of the effective action is $S^{(k)}=\frac12\langle A,\Delta^{(k)}A\rangle+$ the truncation to order $\bar n$ of the formal series $\mathfrak E_k$ of \cite[(3.35)]{Ba1}, built with the contour convention \cite[(3.34)]{Ba1}. This is an algebraic identity between formal power series, proved here. The rescalings \cite[(1.22)]{Ba1} are changes of variables with field-independent Jacobians and change nothing below except additive constants.

\begin{definition}[grading, truncation, re-weighting]\label{def:a06d:B1}
\textup{(a)} In \cite[(3.35)]{Ba1} replace $e'$ by a formal variable $s$ wherever it occurs (on every fluctuation field, and in $\dm(e',\lambda')$, $E(e',\lambda')$), and read \cite[(3.34)]{Ba1} with $j$ from $0$ (which is \cite[(1.3)]{Ba3}). Write $\mathfrak E_k(s,\lambda';A,\phi)\defeq E_k(s,\lambda',eA^{(k),\eps},\phi)$, a formal power series in $(s,\lambda')$ whose coefficients are smooth functions of the block fields; the degree of a monomial is the power of $s$ plus the power of $\lambda'$. \textup{(b)} $T_{\bar n}$ truncates to degree $\le\bar n$ and ``$\mathrm{ev}$'' evaluates at $(s,\lambda')=(e,\lambda)$; in this notation \cite[(3.36)]{Ba1} reads $S^{(k)}=\frac12\langle A,\Delta^{(k)}A\rangle+\mathrm{ev}\,T_{\bar n}\mathfrak E_k$. \textup{(c)} In step $k+1$, $A=A'+aL^{-2}C^{(k)}Q^*B$ \cite[(3.41)]{Ba1} and $A^{(k),\eps}=A'^{(k),\eps}+A^{(k+1),\eps}$ \cite[(3.42)]{Ba1}; for a coefficient function $f$ of the background $eA^{(k),\eps}$, $\iota_kf$ is the Taylor series of $\tau\mapsto f(\tau A'^{(k),\eps}+eA^{(k+1),\eps})$ at $\tau=0$ with $\tau$ then replaced by $s$, so that each power of the new fluctuation carries one power of $s$ (Ba{\l}aban's order count: every $A'$- or $\tilde A$-leg in \cite[(1.8)--(1.15)]{Ba3} carries one factor $e(L^k\eps)$).
\end{definition}

\begin{lemma}[truncation commutes with a Gaussian integral]\label{lem:a06d:B2}
Let $d\mu\propto e^{-F_0}dX$ be a Gaussian law on a finite-dimensional space ($F_0$ a positive quadratic form), and let $W$ be a formal series of degree $\ge1$ whose coefficients are polynomials in $X$. Then $-\log\langle e^{-W}\rangle_\mu$ is a well-defined formal series; its degree-$n$ part is a finite combination of Gaussian moments of the parts $W_1,\dots,W_n$ of $W$ of degrees $1,\dots,n$; hence $T_{\bar n}(-\log\langle e^{-W}\rangle_\mu)=T_{\bar n}(-\log\langle e^{-T_{\bar n}W}\rangle_\mu)$.
\end{lemma}

\begin{proof}
$\langle e^{-W}\rangle=\sum_{m\ge0}(-1)^m\langle W^m\rangle/m!$; the degree-$n$ part of $W^m$ vanishes for $m>n$ and for $m\le n$ involves only $W_1,\dots,W_{n-m+1}$. So $\langle e^{-W}\rangle=1+u$ with $u$ of degree $\ge1$, and the degree-$n$ part of $-\log(1+u)=\sum_{m\ge1}(-1)^mu^m/m$ involves only the parts of $u$ of degree $\le n$. Each coefficient is a Gaussian moment of a polynomial, hence finite.
\end{proof}

This is the cumulant expansion: Ba{\l}aban's ``terms of the order $\le\bar n$'' of an integrated action are exactly $T_{\bar n}$ of this series.

\begin{lemma}[composition of the scalar kernels]\label{lem:a06d:B3}
Let $Q\colon\R^n\to\R^m$ be linear with $QQ^*=c\,\mathbf 1_m$ ($c>0$), and $\alpha,\beta>0$. Then for all $\psi\in\R^m$, $\chi\in\R^n$,
\[
\int_{\R^n}d\phi\,\exp\big(-\tfrac12\alpha|\psi-Q\phi|^2-\tfrac12\beta|\phi-\chi|^2\big)=(2\pi/\beta)^{n/2}(1+\alpha c/\beta)^{-m/2}\exp\big(-\tfrac12\gamma|\psi-Q\chi|^2\big),\qquad\gamma\defeq\frac{\alpha\beta}{\beta+\alpha c}.
\]
\end{lemma}

\begin{proof}
Write $\phi=\chi+\beta^{-1/2}\xi$ with $\xi$ standard Gaussian; then $Q\phi=Q\chi+Z$ with $Z\sim N(0,(c/\beta)\mathbf 1_m)$, and for $Z\sim N(0,\sigma^2\mathbf 1_m)$, $\E\exp(-\frac12\alpha|v+Z|^2)=(1+\alpha\sigma^2)^{-m/2}\exp(-\frac12\frac{\alpha}{1+\alpha\sigma^2}|v|^2)$. Take $v=Q\chi-\psi$, $\sigma^2=c/\beta$.
\end{proof}

The covariant one-step average $(Q(\beta)\phi)(z)=L^{-d}\sum_{y\in B(z)}U(\beta(\Gamma_{z,y}))\phi(y)$ has $Q(\beta)Q(\beta)^*=L^{-d}\mathbf 1$ for every background $\beta$ (the $U$ are orthogonal, the blocks disjoint). With $\alpha=aL^{d-2}$ \cite[(3.38)]{Ba1}, $\beta=a_k$ and $c=L^{-d}$ the lemma gives $\gamma=L^{d-2}aa_k/(a_k+aL^{-2})=L^{d-2}a_{k+1}$, the recursion \cite[(2.13)]{Ba1}; the factor $L^{d-2}$ is removed by the rescaling of $\psi$, and the prefactor depends neither on the background nor on $\psi,\chi$. At a fixed external field this is \cite[(2.10)--(2.16)]{Ba1}; the next two lemmas extend it to level-dependent fields.

\begin{lemma}[composite averages]\label{lem:a06d:B4}
Let $\beta$ be a fine field and $\Phi$ assign a phase to each composite contour $\Gamma^{(k)}_{y,x}$. Then $(Q(\beta)Q_k(\Phi)\phi')(z)=(Q_{k+1}(\Phi')\phi')(z)$ with $\Phi'(\Gamma^{(k+1)}_{z,x})\defeq\beta(\Gamma_{z,x_k})+\Phi(\Gamma^{(k)}_{x_k,x})$.
\end{lemma}

\begin{proof}
$U(a)U(b)=U(a+b)$ since all $U$ commute; $B(z)\circ B^k(y)=B^{k+1}(z)$ and $\Gamma^{(k+1)}_{z,x}=\Gamma_{z,x_k}\cup\Gamma^{(k)}_{x_k,x}$ by \cite[(2.2)]{Ba1}; the weights $L^{-d}\eta^d$ are those of the $(k+1)$-average.
\end{proof}

\begin{lemma}[the contour convention is the iterate, and $j$ starts at $0$]\label{lem:a06d:B5}
For $k\ge1$ put $\Phi_k(s)(\Gamma^{(k)}_{y,x})\defeq eA^{(k),\eps}(\Gamma^{(k)}_{y,x})+s\sum_{j=0}^{k-1}A'^{(j),\eps}(\Gamma^{(j+1)}_{x_{j+1},x})$, which is \cite[(3.34)]{Ba1} with $j$ from $0$ and the fluctuations weighted by $s$. \textup{(a)} With $\beta\defeq\iota_k(eA^{(k),\eps})=sA'^{(k),\eps}+eA^{(k+1),\eps}$ on the first-step contour $\Gamma_{z,x_k}$ prescribed by \cite[(3.37)]{Ba1}, $\beta(\Gamma_{z,x_k})+(\iota_k\Phi_k(s))(\Gamma^{(k)}_{x_k,x})=\Phi_{k+1}(s)(\Gamma^{(k+1)}_{z,x})$. \textup{(b)} At $k=0$ the first kernel carries $U(eA(\Gamma_{z,x}))$ with the full fine field $A=A'_0+A^{(1),\eps}$ \cite[(3.15), p.~615]{Ba1}, and $\iota_0(eA(\Gamma_{z,x}))=sA'_0(\Gamma^{(1)}_{x_1,x})+eA^{(1),\eps}(\Gamma^{(1)}_{z,x})=\Phi_1(s)(\Gamma^{(1)}_{z,x})$: its $s$-term is the $j=0$ term. \textup{(c)} Hence the printed range $j=1,\dots,k-1$ of \cite[(3.34)]{Ba1} is not what the procedure produces: at $k=1$ it would drop $sA'_0(\Gamma^{(1)})$ and with it every contour vertex \cite[(1.14)]{Ba3} with $A'_0$-legs, which the first step does produce \cite[(3.15)--(3.16)]{Ba1}; with $j=0,\dots,k-1$ the convention is the iterate.
\end{lemma}

\begin{proof}
(a) By \cite[(2.2)]{Ba1}, $\Gamma^{(k+1)}_{z,x}=\Gamma_{z,x_k}\cup\Gamma^{(k)}_{x_k,x}$ and $\Gamma^{(k+1)}_{x_{k+1},x}=\Gamma^{(k+1)}_{z,x}$ since $x_{k+1}=z$; by \cite[(2.3)]{Ba1} line integrals add over unions of contours, so $\beta(\Gamma_{z,x_k})+sA'^{(k)}(\Gamma^{(k)}_{x_k,x})+eA^{(k+1)}(\Gamma^{(k)}_{x_k,x})=sA'^{(k)}(\Gamma^{(k+1)}_{x_{k+1},x})+eA^{(k+1)}(\Gamma^{(k+1)}_{z,x})$, and the terms with $j<k$ are unchanged because the points $x_j$, $j\le k$, of $x$ are the same at both levels. (b) and (c) are read off the definitions.
\end{proof}

\begin{lemma}[{the Gaussian chain on a domain; \cite[(3.41), (3.42)]{Ba1} and \cite[(1.27)]{Ba3}}]\label{lem:a06d:B6}
Let $H>0$ be the covariant Neumann operator $-\Delta^\eta_{B,\Omega}+m^2$ on a union $\Omega$ of blocks at a fixed background $B$, or the vector operator at $B=0$ componentwise; let $R_k$ be the $k$-step composite average onto $T^{(k)}_1\cap\Omega$, $Q$ the one-step average with $QQ^*=c\,\mathbf 1$, and $\beta_k,\alpha>0$. Consider the Gaussian density of $(\phi',\phi,\psi)$ proportional to $\exp(-\frac12\langle\phi',H\phi'\rangle-\frac12\beta_k|\phi-R_k\phi'|^2-\frac12\alpha|\psi-Q\phi|^2)$, and put $G_k\defeq(H+\beta_kR_k^*R_k)^{-1}$, $\Delta^{(k)}\defeq\beta_k-\beta_k^2R_kG_kR_k^*$, $C^{(k)}\defeq(\Delta^{(k)}+\alpha Q^*Q)^{-1}$, $\gamma\defeq\alpha\beta_k/(\beta_k+\alpha c)$, $G_{k+1}\defeq(H+\gamma(QR_k)^*(QR_k))^{-1}$. Then: \textup{(a)} given $\phi$, $\phi'\sim N(\beta_kG_kR_k^*\phi,G_k)$, and the marginal density of $\phi$ is $\propto\exp(-\frac12\langle\phi,\Delta^{(k)}\phi\rangle)$; \textup{(b)} given $\psi$, $\phi\sim N(\alpha C^{(k)}Q^*\psi,C^{(k)})$, and $\E[\phi'\mid\psi]=\beta_kG_kR_k^*\alpha C^{(k)}Q^*\psi=\gamma G_{k+1}(QR_k)^*\psi$; \textup{(c)} $\Cov(\phi'\mid\psi)=G_{k+1}=G_k+\beta_k^2G_kR_k^*C^{(k)}R_kG_k$.
\end{lemma}

(a) is the definition of $A^{(k)}$, $\phi^{(k)}$ and $\Delta^{(k)}$ \cite[(3.29), (2.21)]{Ba1}; (b) is \cite[(3.41)]{Ba1} (the conditional fluctuation $A'=A-\E[A\mid B]$ is independent of $B$ with covariance $C^{(k)}$) together with the second equality of \cite[(3.42)]{Ba1}, $a_kG_kQ_k^*\,aL^{-2}C^{(k)}Q^*B=a_{k+1}L^{-2}G^\eta_{k+1}Q^*_{k+1}B$; (c) is, on the torus at $B=0$, one step of \cite[(1.27)]{Ba3}, and on $\Omega$ the same identity with Neumann lines.

\begin{proof}
(a): complete the square in $\phi'$. (b): by (a) the $(\phi,\psi)$-density is $\propto\exp(-\frac12\langle\phi,\Delta^{(k)}\phi\rangle-\frac12\alpha|\psi-Q\phi|^2)$; complete the square in $\phi$. The full density factorises as $f(\phi',\phi)g(\phi,\psi)$, so $\phi'$ and $\psi$ are conditionally independent given $\phi$ and $\E[\phi'\mid\psi]=\E[\E[\phi'\mid\phi]\mid\psi]=\beta_kG_kR_k^*\E[\phi\mid\psi]$. Integrating $\phi$ first (\cref{lem:a06d:B3} with $\chi=R_k\phi'$), the $(\phi',\psi)$-density is $\propto\exp(-\frac12\langle\phi',H\phi'\rangle-\frac12\gamma|\psi-QR_k\phi'|^2)$, whence $\E[\phi'\mid\psi]=\gamma G_{k+1}(QR_k)^*\psi$ and $\Cov(\phi'\mid\psi)=G_{k+1}$. (c): the law of total covariance, $\Cov(\phi'\mid\psi)=\E[\Cov(\phi'\mid\phi,\psi)\mid\psi]+\Cov(\E[\phi'\mid\phi,\psi]\mid\psi)=G_k+\beta_k^2G_kR_k^*C^{(k)}R_kG_k$.
\end{proof}

\begin{lemma}[diagonal truncation]\label{lem:a06d:B7}
Let $H(s,\tau,\lambda')$ be a formal series and $h(s,\lambda')\defeq H(s,s,\lambda')$. Then $T_{\bar n}h$ at $(s,\lambda')=(e,\lambda)$ equals $T^{\mathrm{tot}}_{\bar n}H$ at $(s,\tau,\lambda')=(e,e,\lambda)$, where $T^{\mathrm{tot}}_{\bar n}$ truncates the total degree.
\end{lemma}

\begin{proof}
The coefficient of $s^m\lambda'^\beta$ in $h$ is $\sum_{\alpha+\gamma=m}H_{\alpha\gamma\beta}$.
\end{proof}

\begin{theorem}[{\cite[(3.36)]{Ba1} at every $k$}]\label{thm:a06d:B8}
Let $S^{(0)}\defeq S^\eps$ \cite[(1.20)]{Ba3}, and for $k\ge0$ let $S^{(k+1)}_{\mathrm{pol}}$ be the polynomial part, of order $\le\bar n$, that the step $k+1$ of \cite[Sect.~3]{Ba1} produces from $S^{(k)}_{\mathrm{pol}}$: the kernels of \cite[(3.37)]{Ba1} (gauge $T_{a,L}$, scalar $T_{a,L,A^{(k),\eps}}$), the translation \cite[(3.41)]{Ba1}, the expansion of all coefficients in the new fluctuation \cite[(3.44)]{Ba1}, the Gaussian integral and the cumulant expansion keeping the terms of order $\le\bar n$ (the restrictions $\chi$ and the terms of order $>\bar n$ are split off; they are the error terms of \cite[(3.26)]{Ba1} and are not part of this statement). Then for every $k\ge0$,
\[
S^{(k)}_{\mathrm{pol}}=\tfrac12\langle A,\Delta^{(k)}A\rangle+\mathrm{ev}\,T_{\bar n}\mathfrak E_k+c_k,
\]
where $\mathfrak E_k$ is \cite[(3.35)]{Ba1} with \cite[(3.34)]{Ba1} read from $j=0$, and $c_k$ is a constant independent of the fields and of the couplings (a normalisation, entering $E_0$ only). With the printed range $j\ge1$ the identity is false at $k=1$ (\cref{lem:a06d:B5}(c)).
\end{theorem}

\begin{proof}
Induction on $k$. \emph{Base $k=0$.} $\mathfrak E_0$ contains no fluctuation: $\mathfrak E_0=\frac12\langle\phi,(-\Delta_{eA}+m^2)\phi\rangle+\sum\eps^d(\lambda'|\phi|^4+\frac12\dm(s,\lambda')|\phi|^2)+E(s,\lambda')+\mathrm{const}$, and $\frac12\langle A,\Delta^{(0)}A\rangle=\frac12\langle A,(-\Delta+\lmu^2)A\rangle$; $T_{\bar n}$ acts as the identity, since $\dm$ has orders $2\le\alpha+2\beta\le4\le\bar n$ \cite[p.~417]{Ba3} and $E$ has orders $\le\bar n$. So $S^{(0)}=\frac12\langle A,\Delta^{(0)}A\rangle+\mathrm{ev}\,T_{\bar n}\mathfrak E_0$. At $k=0$ the inner kernel is absent ($Q_0=1$): the gauge part of the step is \cref{lem:a06d:B6}(b) with $H=\Delta^{(0)}$, and the scalar part is \cite[(3.35)]{Ba1} at $k=1$ by definition, with $a_1=a$ and the $j=0$ contour term of \cref{lem:a06d:B5}(b).

\emph{Step $k\to k+1$.} Define
\[
\mathrm{Step}_kF(B,\psi)\defeq-\log\Big[N_k\int dA\int d\phi\,\exp\big(-\tfrac12aL^{d-2}|B-QA|^2-\tfrac12aL^{d-2}|\psi-Q(\iota_keA^{(k),\eps})\phi|^2\big)e^{-\iota_kF(A,\phi)}\Big],
\]
followed by the rescaling \cite[(1.22)]{Ba1}.

(i) \emph{Exact identity} for $F=\frac12\langle A,\Delta^{(k)}A\rangle+\mathfrak E_k$, degree by degree (Fubini for Gaussian moments of polynomials). Gauge part: by \cref{lem:a06d:B6}(b),
$\int dA\exp(-\frac12aL^{d-2}|B-QA|^2-\frac12\langle A,\Delta^{(k)}A\rangle)f(A^{(k),\eps})=N\exp(-\frac12\langle B,\Delta^{(k+1)}B\rangle)\int d\mu_{C^{(k)}}(A'_k)f(A'^{(k),\eps}+A^{(k+1),\eps})$, and $\iota_k$ puts the factor $s$ on $A'^{(k),\eps}$. Scalar part: write $e^{-\mathfrak E_k}$ as the nested integral \cite[(3.35)]{Ba1}; the $\phi$-integral is $\int d\phi\exp(-\frac12\alpha|\psi-Q(\beta)\phi|^2-\frac12a_k|\phi-Q_k(\Phi_k)\phi'|^2)=N'\exp(-\frac12a_{k+1}|\psi-Q(\beta)Q_k(\Phi_k)\phi'|^2)$ by \cref{lem:a06d:B3}, and $Q(\beta)Q_k(\iota_k\Phi_k)=Q_{k+1}(\Phi_{k+1})$ by \cref{lem:a06d:B4,lem:a06d:B5}(a). Covariant Laplacian: on every bond $b$, $(\iota_k\Phi_k)_b=sA'^{(k)}_b+eA^{(k+1)}_b+s\sum_{j<k}A'^{(j)}_b=(\Phi_{k+1})_b$. The $\lambda'$-terms and $\dm(s,\lambda')$, $E(s,\lambda')$ are unchanged. Hence $\mathrm{Step}_k(\frac12\langle A,\Delta^{(k)}A\rangle+\mathfrak E_k)=\frac12\langle B,\Delta^{(k+1)}B\rangle+\mathfrak E_{k+1}+c'_k$.

(ii) \emph{Truncation.} The degree-$0$ part of the integrand of $\mathrm{Step}_k$ is Gaussian in $X=(A'_k,\phi)$: its gauge part is given by (i), its scalar part is $\frac12\langle\phi,\Delta^{(k)}(eA^{(k+1),\eps})\phi\rangle+\frac12\alpha|\psi-Q(eA^{(k+1),\eps})\phi|^2$. Since $\iota_k$ only raises degrees, $T_{\bar n}\iota_kT_{\bar n}=T_{\bar n}\iota_k$, and \cref{lem:a06d:B2} gives $T_{\bar n}\mathrm{Step}_k(T_{\bar n}F)=T_{\bar n}\mathrm{Step}_k(F)$.

(iii) \emph{Evaluation.} Ba{\l}aban's step expands in the new fluctuation at the true coupling $e$ and keeps total order $\le\bar n$; by \cref{lem:a06d:B7} with $\tau$ the weight of the new fluctuation ($\iota_k$ is the diagonal $\tau=s$), his truncation is $\mathrm{ev}\,T_{\bar n}$ of the diagonal. Hence $S^{(k+1)}_{\mathrm{pol}}=\mathrm{ev}\,T_{\bar n}\mathrm{Step}_k(S^{(k)}_{\mathrm{pol}})=\frac12\langle B,\Delta^{(k+1)}B\rangle+\mathrm{ev}\,T_{\bar n}\mathfrak E_{k+1}+c_{k+1}$.
\end{proof}

The theorem is the polynomial-part identity that ``close inspection'' asserts, an implication of the printed definitions \cite[(2.1)--(2.3), (3.29), (3.33), (3.35), (3.37)]{Ba1}, \cite[(1.2)]{Ba4} and of \cite[(3.34)]{Ba1} read with $j$ from $0$, i.e.\ of the defining formula \cite[(1.3)]{Ba3}. It does not cover the analytic content of a step (the restrictions, the order-$>\bar n$ remainders, the error terms of \cite[(3.26)]{Ba1}): that is \cite[Props.~3.1--3.2]{Ba1}, which follow from Proposition~1 (Step~3(D)).
